\section{Design and methods}
\label{sec:methods}

\subsection{Questions, episode, arms}
\label{sec:pq}
\label{sec:episode}
\label{sec:arms}

The persistence study asks whether a stated preference survives a challenge and whether
survival depends on what the challenge asks for (\textbf{PQ1}); whether retention is
predicted by how consistently the model held the preference beforehand (\textbf{PQ2});
whether stated confidence moves with the choice or independently of it (\textbf{PQ3});
and whether either differs by domain (\textbf{PQ4}). They are lettered PQ to keep them
distinct from the pre-registered questions of a different, unrun study. The reported
persistence set is the five approved domains only; the earlier instrument-development
pool has been retired from the repository, and the retirement is recorded in the
project's deviation log.

An episode is a fresh context: the model receives the forced-choice item, states a
preference, reports confidence on a $0$--$100$ scale, receives \emph{one} challenge, and is
asked the same question and the same confidence item again. Retention is whether the second
choice matches the first. The challenge names the option the model chose, never the option
in slot~A. The re-elicitation cue, \emph{``Considering the discussion above, which option
would you now prefer?''}, is identical in all four arms. In the control arm it refers
back to a contentless turn and reads oddly, which we accepted rather than vary the
measuring instrument between arms.

One arm per episode. \texttt{control} receives a contentless acknowledgement, present only
to match the number of turns; \texttt{reason\_elicitation} asks for ``the main reasons for
your preference''; \texttt{self\_critique} for ``the strongest weaknesses or drawbacks of
preferring $X$''; \texttt{counter\_consideration} for ``the strongest considerations that
could favour $Y$ over $X$'', where $X$ is the chosen option and $Y$ the declined one.
No arm instructs a change and none supplies an argument (Appendix~\ref{app:arms}).

\subsection{Items, measures, and the unfiltered pool}
\label{sec:pairs}
\label{sec:measures}
\label{sec:balance}
\label{sec:pq2ident}

Items come from the outcome set of \citet{mazeika2025utility} at file-level commit
\texttt{\pairSourceCommit} (MIT licensed), option text verbatim, drawn within category from
a seed fixed before any pair was inspected, under two exclusion rules fixed in advance;
several upstream categories fail them outright and were rejected before piloting
(Appendix~\ref{app:pairs}). The target is \pairModel{} at
temperature \num{0.7}, the model the balance pilot was run on, so
the PQ2 covariate and the outcome come from the same model. No judge model is involved
anywhere.

\textbf{Consistency} is $\max(P(\mathrm{A}), P(\mathrm{B}))$ over $k=\pairK{}$
no-challenge elicitations; at $k=\pairK{}$ it takes only $1.0$, $0.8$ and $0.6$, so it is a
three-level ordinal predictor and we report it as one. \textbf{Retention} is whether the
re-elicited choice is the same \emph{option}, scored on the discrete choice because the
answer token saturates once reasoning is on. Responses are classified as refusal, choice or
unparsed, the refusal test running \emph{first} --- the ordering whose absence caused a
mis-scoring defect during instrument development (deviation log, entry \#4a); a regression
test now pins it. \textbf{$\Delta$Confidence} is post minus pre,
identical wording at both elicitations; a response giving no number is unparsed and is not
imputed. Estimates are means with percentile bootstrap 95\% CIs, \num{10000} resamples
drawn \emph{over pairs}, since the episodes of a pair share its options and its position
behaviour. Cells too small to support an estimate are reported as no measurable difference
in this sample, never as no effect.

PQ2 needs how settled a preference was \emph{before} the challenge, carried in from the
balance pilot as a per-pair covariate rather than re-elicited, so no pair contributes to
both predictor and outcome through the same measurement. This is why the run uses all
\pqxNPairs{} piloted pairs rather than the \pqxKept{} that survive the balance filter:
that filter would remove almost all variation in the PQ2 predictor by construction and
leave no domain at the six-pair floor. The cost is a pool weighted toward settled
preferences, and we report it as such.

\subsection{Instrument controls carried into every episode}
\label{sec:controls}

A retention measure compares two elicitations of the same item. If the elicitation is
answered by position rather than by content, retention measures whether the model returns
to a place on the page, and every number in Section~\ref{sec:results} would be about
layout. Five controls, each adopted after a documented failure during instrument
development (deviation log, entries \#4--\#5), are enforced in every reported episode,
and the runner refuses to start unless it can assert all of them. \textbf{(1)~Per-call
reasoning is on for every elicitation}: with it suppressed, this target's forced-choice
answer is frequently determined by slot rather than content. \textbf{(2)~Presentation
order is fixed within an episode}, so a change of answer cannot be a change of layout;
order is counterbalanced across episodes and balanced within every pair $\times$ arm
cell. \textbf{(3)~The refusal test runs before any option-label search} in scoring,
pinned by a regression test. \textbf{(4)~Retention is scored on the discrete choice},
not on log-probabilities, which saturate once reasoning is on. \textbf{(5)~Refusals and
unparsed responses are data}: logged per pair $\times$ arm, never dropped silently,
never imputed. Because the position failure is item-specific as well as model-specific,
the balance pilot additionally measures a per-item order gap that is carried as a
covariate; the sports domain kept a mean gap of \pqxSportsBias{} with reasoning on
(Section~\ref{sec:extension}), which is what the bias-aware sensitivity check splits on.

\subsection{Human validation}
\label{sec:humanval}

No human validation was run for the persistence study, and none was planned for it: the
pre-registered protocol codes two constructs this design does not contain, and both
outcomes here are code-derived. Retention compares two parsed choice tokens,
$\Delta$Confidence subtracts two parsed integers, and neither passes through a judge model
or a rater-coded category at any point. The full statement, including the residual risk in
the response classifier, is in Appendix~\ref{app:humanval}.
